\documentclass{reportkit}

\setreportkitleftheader{REPORTKIT / ALGORITHM-VIZ PHASE 1 TEST}
\setreportkitfooter{Algorithm-state visualization phase 1 acceptance test}
\setreportkitversion{v1.10.0}

\title{ReportKit -- algorithm-state visualization, phase 1}
\author{Test build}
\date{16 September 2026}

\begin{document}
\maketitle

\begin{execsummary}
This is an acceptance test for \texttt{reportkit-algorithm-viz.sty}'s Phase 1
milestone: the shared algorithm-state vocabulary, theme tokens,
\texttt{arraystate}, and \texttt{algorithmtrace}. It exercises every public
command in the module and the nine shared states in a single document.
\end{execsummary}

\section{Two-pointer scan}
The motivating example for \texttt{arraystate}: values, two labelled
pointers, and a highlighted active range.

\begin{diagram}[
  label={fig:two-pointer},
  caption={A sorted array with left and right pointers bounding the active range.},
  source={Conceptual diagram.},
  description={A six-element sorted array with the left pointer on index two and the right pointer on index five. The range from index two to five is marked active.}
]
\begin{arraystate}[indices=true]
  \cell{-4}
  \cell{-1}
  \cell{-1}
  \cell{0}
  \cell{1}
  \cell{2}
  \pointer[below]{L}{2}
  \pointer[below]{R}{5}
  \range[state=active]{2}{5}
\end{arraystate}
\end{diagram}

\section{All nine shared states}
Every algorithm-state primitive understands the same nine-state vocabulary.
Meaning is carried by border weight, line style, and shape, not only by
fill colour, so it survives grayscale printing.

\begin{diagram}[
  caption={The nine shared algorithm states: current, active, candidate, frontier, visited, resolved, discarded, blocked, unseen.},
  source={Conceptual diagram.},
  description={Nine cells, one per shared algorithm state, each with a distinct border weight, line style, or shape treatment so the state remains legible in grayscale.}
]
\begin{arraystate}
  \cell[state=current]{Cu}
  \cell[state=active]{Ac}
  \cell[state=candidate]{Ca}
  \cell[state=frontier]{Fr}
  \cell[state=visited]{Vi}
  \cell[state=resolved]{Re}
  \cell[state=discarded]{Di}
  \cell[state=blocked]{Bl}
  \cell[state=unseen]{Un}
\end{arraystate}
\end{diagram}

\section{Prefix sums}
Multi-row mode aligns an original array with a derived row -- here, prefix
sums -- using \texttt{\string\row} instead of \texttt{\string\cell}.

\begin{diagram}[
  label={fig:prefix-sum},
  caption={An array and its prefix sums, with the summed range highlighted.},
  source={Conceptual diagram.},
  description={Two aligned rows: four values and their five-element prefix sums, with the active range on the values row spanning indices one to three.}
]
\begin{arraystate}[rows=2]
  \row{values}{3,1,4,2}
  \row{prefix}{0,3,4,8,10}
  \range[row=values,state=active]{1}{3}
  \annotation{$P_4 - P_1 = 7$}
\end{arraystate}
\end{diagram}

\section{Sliding window as an ordered trace}
\texttt{algorithmtrace} lays out titled snapshots in reading order with a
connecting arrow -- the second-highest-priority Phase 1 primitive after
\texttt{arraystate}.

\begin{diagram}[
  label={fig:sliding-window},
  caption={A fixed-size sliding window advances by one element per step.},
  source={Conceptual diagram.},
  description={Three snapshots show a two-element window advancing one position to the right at each step over a four-element array.}
]
\begin{algorithmtrace}[columns=3]
  \snapshot{Initial}{%
    \begin{arraystate}
      \cell{4}\cell{2}\cell{7}\cell{1}
      \range[state=active]{0}{1}
    \end{arraystate}%
  }
  \snapshot{Advance right}{%
    \begin{arraystate}
      \cell{4}\cell{2}\cell{7}\cell{1}
      \range[state=active]{1}{2}
    \end{arraystate}%
  }
  \snapshot{Advance again}{%
    \begin{arraystate}
      \cell{4}\cell{2}\cell{7}\cell{1}
      \range[state=active]{2}{3}
    \end{arraystate}%
  }
\end{algorithmtrace}
\end{diagram}

\section{Distinguishing the algorithm primitives}
\begin{principle}{An algorithm visual is not an algorithm's control flow}
\texttt{algorithmblock} explains \emph{how} an algorithm decides what to do
next; \texttt{arraystate}/\texttt{algorithmtrace} show \emph{what the data
looks like} while it runs. The two are complementary, not substitutes for
one another.
\end{principle}

\begin{algorithmblock}[label={alg:two-pointer-viz}]{Two-pointer elimination}
  \AlgorithmInput{Heights $h_0,\ldots,h_{n-1}$}
  \AlgorithmOutput{Maximum container area}
  \State $L \gets 0$
  \State $R \gets n-1$
  \While{$L < R$}
    \State $best \gets \max(best,(R-L)\min(h_L,h_R))$
    \If{$h_L \le h_R$}
      \State $L \gets L+1$
    \Else
      \State $R \gets R-1$
    \EndIf
  \EndWhile
  \State \Return $best$
\end{algorithmblock}

\end{document}
